%======================================================================
\section{Why quantum?}\label{sec:motivation}
%======================================================================

Around 1900, physics looked nearly finished. Newton's mechanics described
planets and projectiles, Maxwell's equations described light and electricity,
and thermodynamics described heat. A handful of experiments refused to fit.
Three of them, taken together, are enough to see why a new theory was needed,
and what kind of mathematics it had to use.

\subsection{Three experiments}

\paragraph{The double slit.}
Send a stream of particles towards a wall with two narrow slits, and record
where they land on a screen behind it. If the particles were little bullets,
each one would go through one slit or the other, and the pattern on the screen
would be the sum of two bumps, one behind each slit. Waves behave differently.
A wave goes through \emph{both} slits, the two outgoing waves overlap, and the
screen shows alternating bright and dark bands, called \emph{interference
fringes}, because the waves reinforce each other at some points and cancel at
others (Figure~\ref{fig:slit}).

\begin{figure}[htb]
  \centering
  \begin{minipage}[c]{0.48\textwidth}
    \centering
    \begin{tikzpicture}[scale=0.82]
      % barrier with two slits at y = +-0.5
      \draw[ultra thick,ink] (0,-2.3)--(0,-0.6);
      \draw[ultra thick,ink] (0,-0.4)--(0,0.4);
      \draw[ultra thick,ink] (0,0.6)--(0,2.3);
      % source and incoming rays
      \fill[navy] (-2,0) circle (2.5pt) node[below=3pt,font=\footnotesize] {source};
      \draw[navy,->] (-2,0)--(-0.1,0.5);
      \draw[navy,->] (-2,0)--(-0.1,-0.5);
      % outgoing circular waves from each slit
      \foreach \r in {0.6,1.1,1.6,2.1,2.6,3.1,3.6}{
        \draw[navy!45] ($(0,0.5)+(-50:\r)$) arc (-50:50:\r);
        \draw[navy!45] ($(0,-0.5)+(-50:\r)$) arc (-50:50:\r);
      }
      % screen with fringes drawn as shades
      \foreach \y in {-2.2,-2.15,...,2.15}{
        \pgfmathsetmacro{\o}{0.5*(1+cos(\y*360/1.1))}
        \fill[navy,opacity=\o] (4.05,\y) rectangle (4.5,{\y+0.05});
      }
      \draw[ultra thick,ink] (4,-2.3)--(4,2.3);
      \node[font=\footnotesize] at (4.3,-2.65) {screen};
    \end{tikzpicture}
  \end{minipage}\hfill
  \begin{minipage}[c]{0.48\textwidth}
    \centering
    \begin{tikzpicture}
      \begin{axis}[width=\linewidth, height=5.2cm, xmin=-2.2, xmax=2.2, ymin=0, ymax=4.6,
          xlabel={position $x$ on the screen}, ylabel={intensity},
          ytick=\empty, xtick=\empty, axis lines=left, tick style={draw=none},
          legend style={font=\footnotesize, at={(0.5,1.02)}, anchor=south, draw=none},
          label style={font=\footnotesize}]
        \addplot[navy, very thick, samples=200, domain=-2.2:2.2] {2*(1+cos(deg(2*pi*x/1.1)))};
        \addplot[figacc, very thick, dashed, domain=-2.2:2.2] {2};
        \legend{waves: $\abs{\psi_1+\psi_2}^2$, bullets: $\abs{\psi_1}^2+\abs{\psi_2}^2$}
      \end{axis}
    \end{tikzpicture}
  \end{minipage}
  \caption{Left: the double-slit set-up, with the fringes drawn on the screen. Right: the
  pattern that waves produce (solid) against the pattern that bullets would produce (dashed).}
  \label{fig:slit}
\end{figure}

The surprise is that electrons, photons and even whole molecules produce the
\emph{fringes}, and they do so even when they are sent through the apparatus
one at a time. Each particle arrives at a single point, like a bullet, but
after thousands of arrivals the points build up the interference pattern, as if
each particle had passed through both slits and interfered with itself.

Mathematically this is easy to express if we give each route a complex number.
Let $\psi_1(x)=a\,e^{\ii\varphi_1(x)}$ and $\psi_2(x)=a\,e^{\ii\varphi_2(x)}$ be
the \term{amplitudes} for reaching the point $x$ through slit~1 and slit~2. The
rule is that amplitudes \emph{add}, and the intensity is the squared modulus of
the total:
\[
  \abs{\psi_1+\psi_2}^2=2a^2\bigl(1+\cos(\varphi_1-\varphi_2)\bigr),
\]
which oscillates between $0$ and $4a^2$ as the phase difference changes. This
is precisely the computation of Example~\ref{ex:relative-phase} in the next section: complex numbers
and their phases are where the fringes come from. Bullets would give
$\abs{\psi_1}^2+\abs{\psi_2}^2=2a^2$, the flat dashed line.

\paragraph{Black-body radiation.}
A hot object glows: first red, then orange, then white. A perfect absorber
and emitter of radiation, called a \emph{black body}, glows with a
characteristic spectrum that depends only on its temperature. Classical physics
predicts that the energy emitted at frequency $\nu$ grows like $\nu^2$, so the
total energy would be infinite, the so-called \emph{ultraviolet catastrophe}.
Real bodies do not do this: their curves rise, peak, and fall again.

In 1900 Planck found a formula that matches the measurements exactly, and the
price of the formula was a strange hypothesis: light of frequency $\nu$ can
only be exchanged with matter in whole units of energy $E=h\nu$. With this
assumption the energy at frequency $\nu$ is proportional to
$\nu^3/(e^{h\nu/k_BT}-1)$, which agrees with the classical $\nu^2$ law at low
frequencies and is cut off at high ones (Figure~\ref{fig:blackbody}).

\begin{figure}[htb]
  \centering
  \begin{tikzpicture}
    \begin{axis}[width=0.8\linewidth, height=5.4cm, xmin=0, xmax=8, ymin=0, ymax=3,
        xlabel={frequency, in units of $k_BT/h$}, ylabel={emitted energy},
        ytick=\empty, xtick={0,2,4,6,8}, axis lines=left,
        legend style={font=\footnotesize, at={(0.04,0.96)}, anchor=north west, draw=none},
        label style={font=\footnotesize}, tick label style={font=\footnotesize}]
      \addplot[figacc, very thick, dashed, domain=0:1.74, samples=60] {x^2};
      \addplot[navy, very thick, samples=120, domain=0.01:8] {x^3/(exp(x)-1)};
      \legend{classical, Planck}
    \end{axis}
  \end{tikzpicture}
  \caption{Spectrum of a black body. The classical prediction (dashed, cut off at the plot's top)
  keeps growing; Planck's formula (solid) peaks and falls. The curves agree at low frequency.}
  \label{fig:blackbody}
\end{figure}

\paragraph{The stability of the atom.}
In Rutherford's picture of the hydrogen atom, an electron orbits a nucleus. But
an orbiting charge is an accelerating charge, and Maxwell's equations say that
accelerating charges radiate. Classically the electron would lose energy and
spiral into the nucleus in about $10^{-11}$ seconds: atoms could not exist.
Moreover, hydrogen gas does not glow with a continuous spectrum but with a
handful of sharp lines, at frequencies that fit the pattern
$E_n=-13.6\,\mathrm{eV}/n^2$ for the energy levels of the atom
(Figure~\ref{fig:levels}): only certain energies are allowed, and light is
emitted exactly when the electron jumps between them.

\begin{figure}[htb]
  \centering
  \begin{tikzpicture}[scale=1.0]
    % vertical positions are not to scale: n=1 -> -4.6, n=2 -> -1.6, n=3 -> -0.95, n=4 -> -0.55
    \foreach \n/\E/\y in {1/-13.6/-4.6,2/-3.4/-1.6,3/-1.51/-0.95,4/-0.85/-0.55}{
      \draw[ink, thick] (0,\y) -- (4.2,\y);
      \node[left,font=\footnotesize] at (0,\y) {$n=\n$};
      \node[right,font=\footnotesize] at (4.2,\y) {$\E$\,eV};
    }
    \draw[ink, thick, dotted] (0,0) -- (4.2,0);
    \node[right,font=\footnotesize] at (4.2,0) {$0$};
    \draw[->,navy,thick] (0.7,-1.6) -- (0.7,-4.6);
    \draw[->,navy,thick] (1.3,-0.95) -- (1.3,-4.6);
    \draw[->,navy,thick] (1.9,-0.55) -- (1.9,-4.6);
    \node[font=\footnotesize,navy,align=center] at (1.3,-5.2) {ultraviolet lines\\(towards $n=1$)};
    \draw[->,figacc,thick] (2.9,-0.95) -- (2.9,-1.6);
    \draw[->,figacc,thick] (3.5,-0.55) -- (3.5,-1.6);
    \node[font=\footnotesize,figacc,align=center] at (3.2,-2.4) {visible lines\\(towards $n=2$)};
    \node[font=\scriptsize,ink] at (2.1,0.35) {(heights not to scale)};
  \end{tikzpicture}
  \caption{Energy levels of hydrogen. Only the energies shown are allowed. Light is emitted
  when the electron drops from one level to a lower one; the drops towards $n=2$ give the
  visible lines.}
  \label{fig:levels}
\end{figure}

Each experiment has the same moral. Energy comes in units, behaviour is
described by amplitudes that add, and what we observe are squared moduli of
those amplitudes. We need a theory whose basic objects are complex-valued
\emph{waves of probability amplitude}, and whose basic rule is addition.

\subsection{Schr\"odinger's equation}

In 1926 Schr\"odinger wrote down the equation governing such waves for a
particle of mass $m$ moving in a potential $V$:
\begin{equation}\label{eq:schrodinger}
  \ii\hbar\,\frac{\partial\psi}{\partial t}(x,t)
  =\hat H\psi(x,t),\qquad
  \hat H=-\frac{\hbar^2}{2m}\frac{\partial^2}{\partial x^2}+V(x).
\end{equation}
The unknown $\psi(x,t)$ is complex-valued, and it is called the \term{wave
function}. According to Born, $\abs{\psi(x,t)}^2$ is the probability density of
finding the particle near $x$ at time $t$. We take this as motivation here; in
Chapter~\ref{ch:postulates} it becomes a postulate.

Let us look at the simplest case, a free particle ($V=0$).

\begin{example}\label{ex:planewave}
  Show that $\psi(x,t)=A\,e^{\ii(kx-\omega t)}$ with $A\in\C$ solves the free
  equation if $\hbar\omega=\hbar^2k^2/2m$.
\end{example}
\begin{solution}
  Differentiate: $\partial_t\psi=-\ii\omega\psi$ and $\partial_x^2\psi=-k^2\psi$. Substituting
  into $\ii\hbar\,\partial_t\psi=-\frac{\hbar^2}{2m}\partial_x^2\psi$,
  \[
    \ii\hbar(-\ii\omega)\psi=\hbar\omega\,\psi
    \qquad\text{and}\qquad
    -\frac{\hbar^2}{2m}(-k^2)\psi=\frac{\hbar^2k^2}{2m}\,\psi,
  \]
  so the equation holds exactly when $\hbar\omega=\hbar^2k^2/2m$. Note also that
  $\abs{\psi}^2=\abs{A}^2$ does not depend on $x$ or $t$.
\end{solution}

The relation between $\omega$ and $k$ is the free particle's energy
$E=\hbar\omega=p^2/2m$ with momentum $p=\hbar k$. That a plane wave has a
constant $\abs{\psi}^2$ also reveals a difficulty: it cannot be normalised
over all of $\R$, so it is an idealisation.

Now the observation that drives this whole chapter. Equation
\eqref{eq:schrodinger} is \emph{linear} in $\psi$.

\begin{example}\label{ex:linear-schrodinger}
  Show that if $\psi_1$ and $\psi_2$ solve \eqref{eq:schrodinger}, so does
  $\alpha\psi_1+\beta\psi_2$ for all $\alpha,\beta\in\C$. Does the product $\psi_1\psi_2$ solve it?
\end{example}
\begin{solution}
  Both $\partial_t$ and $\hat H$ are linear operations, so
  \[
    \ii\hbar\,\partial_t(\alpha\psi_1+\beta\psi_2)
    =\alpha\,\ii\hbar\,\partial_t\psi_1+\beta\,\ii\hbar\,\partial_t\psi_2
    =\alpha\hat H\psi_1+\beta\hat H\psi_2
    =\hat H(\alpha\psi_1+\beta\psi_2).
  \]
  The product does not work in general: $\hat H$ contains a second derivative,
  and $(\psi_1\psi_2)''=\psi_1''\psi_2+2\psi_1'\psi_2'+\psi_1\psi_2''$ produces a cross term
  that has no counterpart on the other side.
\end{solution}

So the set of solutions is closed under sums and complex multiples: it is a
\emph{complex vector space}. That is the same statement as the double slit, in
which the amplitude for ``both slits'' is the \emph{sum} of the amplitudes for
each. Superposition is not a mysterious extra; it is linearity.

Finally, one special family of solutions will matter a great deal. If $\hat H$
does not depend on time, look for solutions of the form
$\psi(x,t)=\varphi(x)\,T(t)$. Substituting, one finds that $\varphi$ must satisfy
$\hat H\varphi=E\varphi$ for a constant $E$, and that $T(t)=e^{-\ii Et/\hbar}$
(Exercise~\ref{ex:separation}). Thus
\[
  \psi(x,t)=\varphi(x)\,e^{-\ii Et/\hbar},\qquad \hat H\varphi=E\varphi .
\]
The equation $\hat H\varphi=E\varphi$ is an \emph{eigenvalue problem}: $\varphi$ is an eigenvector of
the operator $\hat H$, and $E$, the energy, is its eigenvalue. The discrete
levels of Figure~\ref{fig:levels} are exactly the eigenvalues of the atom's
Hamiltonian. These \term{stationary states} only get multiplied by a phase as
time goes on, so $\abs{\psi}^2$ does not change.

This is the programme of the chapter. Solutions of a linear equation live in
a vector space; probabilities come from complex numbers through their moduli;
and the questions we can ask about the system are eigenvalue problems. Each of
these three ideas gets its own section. We begin, in the next one, with the
numbers.

\begin{exerciseset}
  \exercise{1} In the double slit, let $\psi_j(x)=a\,e^{\ii\varphi_j(x)}$ with $a>0$. Show that
  $\abs{\psi_1+\psi_2}^2=2a^2\bigl(1+\cos(\varphi_1-\varphi_2)\bigr)$, and find the phase
  differences for which the screen is dark.

  \exercise{1} Let $\psi=A\,e^{\ii(kx-\omega t)}$. Compute $\partial_t\psi$ and $\partial_x^2\psi$,
  substitute into the free equation, and derive $\omega=\hbar k^2/2m$. Is $\psi$ normalisable over $\R$?

  \exercise{2} Let $\hat H=-\frac{\hbar^2}{2m}\partial_x^2+V(x)$. Prove that the set of solutions
  of $\ii\hbar\,\partial_t\psi=\hat H\psi$ is a complex vector space. Give a concrete example of
  two solutions of the free equation whose product is not a solution.

  \exercise{2}\label{ex:separation} Assume $\hat H$ does not depend on $t$ and look for solutions
  $\psi(x,t)=\varphi(x)T(t)$. Show that $\ii\hbar\,T'(t)/T(t)=(\hat H\varphi)(x)/\varphi(x)$, so both
  sides equal a constant $E$. Solve for $T$, and prove that $\abs{\psi}^2$ does not depend on $t$.

  \exercise{2} Let $\hat H\varphi_j=E_j\varphi_j$ for $j=1,2$ with $E_1\neq E_2$ and $\varphi_j$
  real-valued, and let $\psi=c_1\varphi_1e^{-\ii E_1t/\hbar}+c_2\varphi_2e^{-\ii E_2t/\hbar}$ with
  $c_1,c_2>0$. Show that $\psi$ solves the time-dependent equation, and that $\abs{\psi}^2$
  oscillates with angular frequency $(E_2-E_1)/\hbar$. Is $\psi$ a stationary state?
\end{exerciseset}
